\documentclass{article}
\usepackage{amsmath,amssymb,hyperref}
\title{When passive geometric tracks cannot recover a named duty}
\author{}
\date{26 November 2026}
\begin{document}
\maketitle
\paragraph{Audit.}
A held-in R1/R6 discriminator in a deliberately small continuous plant.
There are two indistinguishable discs with first-order speed limit $V$
and minimum common center separation $d>0$. The plane otherwise has
no obstacles; a geometric link bound at least $r$ is allowed. At both
visible cuts the midpoint is the same and the two anonymous exact
detections have separation $r\ge d$. Between those cuts the optical
output is blank, not a stream that reports smooth track velocities or
per-object appearances. The yaw, acceleration, imaging, and clearance
uncertainty are outside this model. There is no RF or previously
exchangeable identity information during the blind interval.
An initial named anchor is supplied before the first gap. The named
act is release of an existing duty by its uniquely authorized owner,
not an anonymous request accepting either disc.

\paragraph{Sharp joint motion bound.}
Fix the coordinate orientation at each visible cut. The two associations
map the relative vector (length $r$) through unsigned central angles
$\alpha$ and $\pi-\alpha$, respectively. Write $\theta\in[0,\pi]$
for either and $\beta=\arccos(d/r)$. The exact minimum time is
$L(\theta)/(2V)$ where
\[
 L(\theta)=
 \begin{cases}
 2r\sin(\theta/2),&\theta\le 2\beta,\\[2pt]
 2\sqrt{r^2-d^2}+d(\theta-2\beta),&\theta>2\beta .
 \end{cases}
\]
The relative vector moves at speed at most $2V$ and cannot enter the
open disc of radius $d$. Its shortest exterior path is the chord when
safe, otherwise two tangents and the intervening shorter arc. To see
the latter, any minimizing path is straight off the disc and cannot
have an interior corner off the disc. Once it first meets the boundary,
shortening any intervening excursion by radial projection gives a
boundary arc. In the nonchord case the entry and exit segments are
tangents: moving the contact point toward the free chord shortens the
segment plus arc until the tangent limit. This also follows by
differentiating each contact length plus remaining arc, whose
derivative vanishes at tangency. The other side has no shorter arc
because $\theta\le\pi$. At $r=d$, use the boundary arc alone.
The bounds are attained for the original pair by using half the
relative path at each side of the fixed midpoint, with opposite signs.
The relative path never exceeds radius $r$, so requiring geometric
link reach at least $r$ does not change the result.

This formula tests coexistence of both tracks, not merely that the
two independent speed cones admit the association. With $r=1$,
$d=1/2$, $V=1$, and the same two anonymous positions on reappearance,
independent matching allows swapped labels at time 1. The physical swap
first becomes possible at about 1.127824792. At 1.08 it is impossible
while the unchanged association is possible. At 1.20 both actual
physical continuations produce the same input data. Choosing the
nearest or lowest-cost assignment is a point estimator, not proof
that only that assignment is possible.

\paragraph{What an anonymous history can and cannot authorize.}
With these memoryless first-order movements and no hidden couplings
between successive gaps, each blind interval contributes a set of
allowed permutations. Updates compose their support by exclusive-or,
not by choosing the most likely track afresh. Once both parities are
supported, a later nonempty anonymous transition, even if its
geometric correspondence is unique, cannot recover the original
absolute name. A credible named observation at one distinguishable
endpoint removes the ambiguity if it arrives at the deciding actor
before use. It is one bit for this two-alternative question, not a
mandate to reveal all track metadata. It cannot retroactively justify
an already executed release. Authenticated owner authorization,
physical settlement and role freshness remain separate; observing a
label does not in itself authorize that duty. A physical interlock at
the actual duty boundary or a proven no-exchange constraint is a
different way to avoid depending on recovered vision identity.
If velocity or a continuing commitment couples gaps, the hypothesis
allowing every concatenation has to be retested.

\paragraph{Comparators and evidence.}
\href{https://arxiv.org/abs/1607.04005}{Kirkpatrick and Liu} characterize
minimum total trace length for two labelled discs. The symmetric
midpoint reduction here makes the same relative length a makespan bound
as well. The geometry is not asserted to be new; it is used to compute
the observation support that actually matters for a named act.
\href{https://www.sciencedirect.com/science/article/pii/S0005109815003490}{Yin and Lafortune}
give the more general dynamic-observation supervisory-control setting;
the present binary support calculation ties a complete direct history
enumerator and is not a new general observer algorithm.
The geometric replay independently brackets path lengths using
inscribed and circumscribed polygon visibility graphs, under explicit
floating-point tolerances, on 105 cases. It compares 2187 anonymous
transport histories to exhaustive parity enumeration. Lean checks the
three logical kernels about information loss and same-view decisions.
This is modeled non-RF evidence, not native vision testing or the
independently selected non-radio case required by R1.
\end{document}
